\documentclass[11pt]{article}
\usepackage[a4paper,margin=22mm]{geometry}
\usepackage[no-math]{fontspec}
\setmainfont{lmroman10-regular.otf}[BoldFont=lmroman10-bold.otf,ItalicFont=lmroman10-italic.otf,BoldItalicFont=lmroman10-bolditalic.otf]
\usepackage{amsmath,amssymb,amsthm,mathtools,graphicx,xcolor,hyperref,xurl,xspace,ifthen}
\setlength{\emergencystretch}{4em}
\SetMathAlphabet{\mathbf}{normal}{OT1}{cmr}{bx}{n}
\SetMathAlphabet{\mathbf}{bold}{OT1}{cmr}{bx}{n}
\newtheorem{theo}{Theorem}[section]
\newtheorem{lemm}[theo]{Lemma}
\newtheorem{theorem}[theo]{Theorem}
\newtheorem*{remas*}{Remarks}
\newtheorem{remas}[theo]{Remarks}
\newtheorem*{exam*}{Example}
\newtheorem{ques}[theo]{Question}
\providecommand{\og}{«\nobreakspace}
\providecommand{\fg}{\nobreakspace»}
\newtheorem{lem}[theo]{Lemma}
\newtheorem{lemma}[theo]{Lemma}
\newtheorem{prop}[theo]{Proposition}
\newtheorem{coro}[theo]{Corollary}
\newtheorem{corol}[theo]{Corollary}
\newtheorem{conj}[theo]{Conjecture}
\newtheorem{Proposition}[theo]{Proposition}
\newtheorem{theoreme}[theo]{Theorem}
\theoremstyle{definition}
\newtheorem{defi}[theo]{Definition}
\newtheorem{exam}[theo]{Example}
\newtheorem{exem}[theo]{Example}
\newtheorem{remark}[theo]{Remark}
\newtheorem{example}[theo]{Example}
\newtheorem{rema}[theo]{Remark}
\newtheorem{nota}[theo]{Notation}
\newtheorem*{rema*}{Remark}
\newtheorem*{defi*}{Definition}
\newtheorem*{theo*}{Theorem}
\newtheorem*{lemm*}{Lemma}
\newtheorem*{prop*}{Proposition}
\newcommand{\enoncename}{}
\newtheorem{corpusenonce}[theo]{\enoncename}
\newenvironment{enonce}[2][]{\renewcommand{\enoncename}{#2}\begin{corpusenonce}}{\end{corpusenonce}}
\newenvironment{enonce*}[2][]{\par\medskip\noindent\textbf{#2.}\itshape}{\par\medskip}
\newenvironment{proofc}[1][\proofname]{\begin{proof}[#1]}{\end{proof}}
\newcommand{\Mc}{\mathcal{M}}
\newcommand{\pointir}{.\ ---\ }
\providecommand{\mathof}[1]{\mathrm{#1}}

\numberwithin{equation}{section}
\numberwithin{figure}{section}
\numberwithin{table}{section}
\def\endappendix{}
%~Mouliné par MaN_auto v.0.21.4 2019-04-30 15:11:50


\usepackage{tikz}
\usepackage{mathtools}

\DeclareMathOperator*{\union}{\bigcup}
\DeclareMathOperator*{\inter}{\bigcap}
\DeclareMathOperator{\var}{\ensuremath{\PEfont Var}}
\DeclareMathOperator{\cov}{\ensuremath{\PEfont Cov}}
\DeclareMathOperator{\sign}{sign}
\DeclareMathOperator{\re}{Re}
\DeclareMathOperator{\im}{Im}

\usetikzlibrary{patterns}
\usetikzlibrary{arrows}








\newcommand{\bbE}{{\ensuremath{\mathbb E}}}
\newcommand{\bbK}{{\ensuremath{\mathbb K}}}
\newcommand{\bbN}{{\ensuremath{\mathbb N}}}
\newcommand{\bbP}{{\ensuremath{\mathbb P}}}
\newcommand{\bbR}{{\ensuremath{\mathbb R}}}
\newcommand{\bbZ}{{\ensuremath{\mathbb Z}}}

\newcommand{\cD}{{\ensuremath{\mathcal D}}}
\newcommand{\cF}{{\ensuremath{\mathcal F}}}
\newcommand{\cG}{{\ensuremath{\mathcal G}}}
\newcommand{\cI}{{\ensuremath{\mathcal I}}}
\newcommand{\cL}{{\ensuremath{\mathcal L}}}
\newcommand{\cP}{{\ensuremath{\mathcal P}}}
\newcommand{\cZ}{{\ensuremath{\mathcal Z}}}

\newcommand{\ga}{\alpha}
\newcommand{\gb}{\beta}
\newcommand{\gga}{\gamma}
\newcommand{\gd}{\delta}
\newcommand{\gD}{\Delta}
\newcommand{\gep}{\varepsilon}
\newcommand{\gz}{\zeta}
\newcommand{\gt}{\theta}
\newcommand{\gl}{\lambda}
\newcommand{\gs}{\sigma}
\newcommand{\go}{\omega}

\renewcommand{\tilde}{\widetilde}
\DeclareMathSymbol{\leqslant}{\mathalpha}{AMSa}{"36}
\DeclareMathSymbol{\geqslant}{\mathalpha}{AMSa}{"3E}
\DeclareMathSymbol{\eset}{\mathalpha}{AMSb}{"3F}

\newcommand{\dd}{\mathrm{d}}
\newcommand*{\dx}{\dd x}
\newcommand*{\dy}{\dd y}

\newcommand{\R}{\mathbb{R}}
\newcommand{\N}{\mathbb{N}}

\newcommand*\bs{\boldsymbol}

\newcommand{\PEfont}{\mathrm}

\newcommand{\p}{\ensuremath{\mathrm{P}}}
\newcommand\bP{\ensuremath{\mathrm{P}}}
\newcommand\bE{\ensuremath{\mathrm{E}}}

\newcommand{\ind}{{\sf 1}}

\newcommand{\ef}{\mathrm{eff}}
\newcommand{\hu}{\hat u}
\newcommand{\hU}{\hat U}

\newcommand{\htau}{\hat\tau}
\newcommand{\hE}{\hat\bE}
\newcommand{\hP}{{\hat\bP}}
\newcommand{\ha}{\hat\ga}
\newcommand{\hK}{\hat K}
%\newcommand{\hL}{\hat L}
\newcommand{\hmu}{\hat \mu}
\newcommand{\gann}{\gamma_{\mathrm{ann}}}
\newcommand{\gap}{{\mathrm{gap}}}
\newcommand{\sm}{{\raise0.3ex\hbox{$\scriptstyle \setminus$}}}

\newcommand*{\II}{{I\mk I}}

\graphicspath{{./figures/}}

\newcommand*{\mk}{\mkern -1mu}
\newcommand*{\Mk}{\mkern -2mu}
\newcommand*{\mK}{\mkern 1mu}
\newcommand*{\MK}{\mkern 2mu}

\hypersetup{urlcolor=purple, linkcolor=blue, citecolor=red}

\newcommand*{\romanenumi}{\renewcommand*{\theenumi}{\roman{enumi}}}
\newcommand*{\Romanenumi}{\renewcommand*{\theenumi}{\Roman{enumi}}}
\newcommand*{\alphenumi}{\renewcommand*{\theenumi}{\alph{enumi}}}
\newcommand*{\Alphenumi}{\renewcommand*{\theenumi}{\Alph{enumi}}}
\let\oldforall\forall
\renewcommand*{\forall}{\mathrel{\oldforall}}

\newenvironment{lemmc}{\renewcommand{\pointir}{}\begin{lemm}}{\end{lemm}}
\newenvironment{step}[1]{\begin{proof}[#1]}{\let\qed\relax \end{proof}}
\newenvironment{laststep}[1]{\begin{proof}[#1]}{\end{proof}}

\makeatletter
\let\c@cdrthm\c@theo
\makeatother
\usepackage{cleveref}
\crefname{originallemma}{original lemma}{original lemmas}
\Crefname{originallemma}{Original lemma}{Original lemmas}
\makeatletter
\newlabel{lem:comp.tilt.shift}{{4.3}{source}{}{sourcepointer-1}{}}
\newlabel{lem:comp.tilt.shift@cref}{{[originallemma][0][]4.3}{[1][0][]source}{}{}{}}
\makeatother
\crefname{originallemma}{original lemma}{original lemmas}
\Crefname{originallemma}{Original lemma}{Original lemmas}
\makeatletter
\newlabel{lem:mon.pro}{{4.2}{source}{}{sourcepointer-2}{}}
\newlabel{lem:mon.pro@cref}{{[originallemma][0][]4.2}{[1][0][]source}{}{}{}}
\makeatother
\crefname{originalproposition}{original proposition}{original propositions}
\Crefname{originalproposition}{Original proposition}{Original propositions}
\makeatletter
\newlabel{pr:exist.que.free.energy}{{2.7}{source}{}{sourcepointer-3}{}}
\newlabel{pr:exist.que.free.energy@cref}{{[originalproposition][0][]2.7}{[1][0][]source}{}{}{}}
\makeatother
\crefname{originalsection}{original section}{original sections}
\Crefname{originalsection}{Original section}{Original sections}
\makeatletter
\newlabel{sec2.2}{{2.2}{source}{}{sourcepointer-4}{}}
\newlabel{sec2.2@cref}{{[originalsection][0][]2.2}{[1][0][]source}{}{}{}}
\makeatother
\crefname{originalsection}{original section}{original sections}
\Crefname{originalsection}{Original section}{Original sections}
\makeatletter
\newlabel{sec4}{{4}{source}{}{sourcepointer-5}{}}
\newlabel{sec4@cref}{{[originalsection][0][]4}{[1][0][]source}{}{}{}}
\makeatother
\crefname{originaltheorem}{original theorem}{original theorems}
\Crefname{originaltheorem}{Original theorem}{Original theorems}
\makeatletter
\newlabel{thm:smoothing}{{2.8}{source}{}{sourcepointer-6}{}}
\newlabel{thm:smoothing@cref}{{[originaltheorem][0][]2.8}{[1][0][]source}{}{}{}}
\makeatother
\title{Annealed critical-curve asymptotics for renewal-correlated pinning}
\author{Dimitris Cheliotis, Yuki Chino, Julien Poisat}
\date{doi:10.5802/ahl.11; CC BY 4.0}
\begin{document}
\maketitle
\noindent\textit{Editorial note.} Only the selected claim is counted. Other author results are retained as context or local proof support; no mathematical content has been rewritten.\par\medskip
\noindent\textit{Allowed standard prerequisites (editorial).}\begin{itemize}
\item Renewal theory, homogeneous pinning free energy, Tauberian estimates, basic Fourier analysis and probability inequalities as cited by the authors. Standard cited foundations only; all selected local constructions and proof estimates remain in source ranges.
\end{itemize}\medskip
\ifdefined\hypertarget\hypertarget{sourcepointer-1}{}\fi
\noindent\textit{Source reference.} Original lemma 4.3: focused-seeds/AGP/sources/probability-continuation-20261011/2687-source.tex, source lines 810--810. Original cross-reference in retained introductory navigation or a declared standard renewal input. The selected claim and every local core estimate used for it are retained in full; this pointer does not substitute for a selected proof.\par\medskip
\ifdefined\hypertarget\hypertarget{sourcepointer-2}{}\fi
\noindent\textit{Source reference.} Original lemma 4.2: focused-seeds/AGP/sources/probability-continuation-20261011/2687-source.tex, source lines 781--781. Original cross-reference in retained introductory navigation or a declared standard renewal input. The selected claim and every local core estimate used for it are retained in full; this pointer does not substitute for a selected proof.\par\medskip
\ifdefined\hypertarget\hypertarget{sourcepointer-3}{}\fi
\noindent\textit{Source reference.} Original proposition 2.7: focused-seeds/AGP/sources/probability-continuation-20261011/2687-source.tex, source lines 285--285. Original cross-reference in retained introductory navigation or a declared standard renewal input. The selected claim and every local core estimate used for it are retained in full; this pointer does not substitute for a selected proof.\par\medskip
\ifdefined\hypertarget\hypertarget{sourcepointer-4}{}\fi
\noindent\textit{Source reference.} Original section 2.2: focused-seeds/AGP/sources/probability-continuation-20261011/2687-source.tex, source lines 281--281. Original cross-reference in retained introductory navigation or a declared standard renewal input. The selected claim and every local core estimate used for it are retained in full; this pointer does not substitute for a selected proof.\par\medskip
\ifdefined\hypertarget\hypertarget{sourcepointer-5}{}\fi
\noindent\textit{Source reference.} Original section 4: focused-seeds/AGP/sources/probability-continuation-20261011/2687-source.tex, source lines 674--674. Original cross-reference in retained introductory navigation or a declared standard renewal input. The selected claim and every local core estimate used for it are retained in full; this pointer does not substitute for a selected proof.\par\medskip
\ifdefined\hypertarget\hypertarget{sourcepointer-6}{}\fi
\noindent\textit{Source reference.} Original theorem 2.8: focused-seeds/AGP/sources/probability-continuation-20261011/2687-source.tex, source lines 291--291. Original cross-reference in retained introductory navigation or a declared standard renewal input. The selected claim and every local core estimate used for it are retained in full; this pointer does not substitute for a selected proof.\par\medskip
\section{Introduction}

The goal of this paper is to study the phase transition of the pinning model in presence of a correlated disorder sequence built out of a renewal sequence. We first present the general set-up of pinning models before introducing our specific model. For a review on pinning models, we refer to the three monographs~\cite{Gi2007, Gi2011, dH09} and references therein. In this paper we write $\bbN = \{1, 2, \ldots\}$ and $\bbN_0 = \{0, 1, 2,\ldots\}$.

\subsection{General set-up}

The pinning model provides a general mathematical framework for studying various physical phenomena such as the wetting transition of interfaces, DNA denaturation, or (de)localization of a polymer along a defect line. This statistical-mechanical model is formulated in terms of a Markov chain $(S_n)_{n\in\bbN_0}$ which is given a reward/penalty $\omega_n$ (depending on the sign) when it returns to its initial state $0$ at time $n$.

Let us denote by $\tau = (\tau_n)_{n\in\bbN_0}$ the sequence of return times to $0$, whose law is denoted by $\bP$. It is a renewal sequence starting at $\tau_0 = 0$, and we assume that the inter-arrival law satisfies
\begin{equation}\label{eq:defK_regvar}
K(n) := \bP(\tau_1 = n) = L(n) \, n^{- (1 + \ga)},\qquad \ga > 0,\quad n\in\bbN,
\end{equation}
where $L$ is a slowly varying function whose support is aperiodic, that is, $\gcd\{n \geq 1 \colon L(n) > 0\}= 1$. We also assume that the renewal process is recurrent, that is $\bP(\tau_1 < \infty) = \sum_{n\ge 1} K(n) = 1$ (otherwise it is said to be transient). By a slight abuse of notation, we shall use $\tau$ to refer to the set $\{\tau_k\}_{k\in\bbN_0}$ and write $\gd_n = \ind_{\{n\in\tau\}}$. Independently of $\tau$, we introduce a disorder sequence, that is a sequence of real valued random variables $\go = (\go_n)_{n\in \bbN_0}$ whose law is denoted by $\bbP$.

The object of interest is the sequence of Gibbs measures, also called polymer measures, defined by:
\begin{equation}\label{eq:defquemeas}
\frac{\dd \bP_{n,\gb,h}}{\dd \bP} = \frac{1}{Z_{n,\gb,h}} \exp \bigg\{\sum_{k=1}^n (h + \gb \go_k)\gd_k \bigg\} \,\gd_n, \quad n \in \bbN, \quad \gb \geq 0, \quad h \in \bbR,
\end{equation}
where
\begin{equation}\label{eq:defquepf}
Z_{n,\gb,h} = \bE \left(e^{\sum_{k=1}^n (h + \gb \go_k) \gd_k} \gd_n \right)
\end{equation}
is the quenched partition function, $h$ is called a pinning strength or chemical potential, and $\gb$ is the inverse temperature.

The free energy of the model is defined by
\begin{equation}
F(\gb,h) = \lim_{n \to \infty} \frac{1}{n} \log Z_{n,\gb,h} \geq 0,
\end{equation}
where the limit holds $\bbP$-a.s.~and in $L^1(\bbP)$ under rather mild assumptions on $\go$, namely if $\go$ is a stationary and ergodic sequence of integrable random variables. Then the two phases of the model are the localized phase $\cL = \{(\gb, h): F(\gb, h)>0\}$, where the contact fraction $\partial_h F(\gb,h) = \lim_{n \to \infty} (1/n) \bE_{n,\gb,h}(|\tau \cap \{1,\ldots,n\}|)$ is positive, and the delocalized phase $\cD = \{(\gb, h): F(\gb, h)=0\}$, where it is zero. The two main features of the transition are the quenched critical point and the critical exponent:
\begin{equation}
h_c(\gb) = \inf\{h \colon F(\gb,h) > 0\},\quad \nu_q(\gb) = \lim_{h \searrow h_c(\gb)} \frac{\log F(\gb,h)}{\log (h - h_c(\gb))},
\end{equation}
when the limit exists. The critical curve separates the two phases whereas the critical exponent indicates how smooth the transition is between them.

\subsubsection{Disorder relevance}

One reason for the success of this model is the solvable nature of the homogeneous case, which corresponds to the choice $\gb = 0$ and which is treated in detail in~\cite{Gi2007}. For the moment, we recall that $h_c(0) = 0$ and $\nu_{\hom}:= \nu_q(0) = \max(1,1/\ga)$, see~\cite[Theorem~2.1]{Gi2007}.

An important challenge in statistical mechanics is to understand the effect of quenched impurities or inhomogeneities in the interaction on the mechanism of the phase transition. This can be done by comparing the critical features of the quenched model to that of the annealed model, which is defined by
\begin{equation}\label{eq:defannmeas}
\frac{\dd \bP^a_{n,\gb,h}}{\dd \bP} = \frac{1}{Z^a_{n,\gb,h}} \bbE\left(\exp \left\{ \sum_{k=1}^n (h + \gb \go_k) \gd_k \right\} \gd_n \right), \quad n \in \bbN, \quad \gb \geq 0,\quad h \in \bbR,
\end{equation}
where
\begin{equation}\label{eq:defannpf}
Z^a_{n,\gb,h} = \bbE\bE\left(e^{\sum_{k=1}^n (h+ \gb \go_k)\gd_k} \gd_n\right) = \bbE (Z_{n, \gb, h})
\end{equation}
is the annealed partition function, and the annealed free energy is also defined by
\begin{equation}
F^a(\gb,h) = \lim_{n \to \infty} \frac{1}{n} \log Z^a_{n,\gb,h},
\end{equation}
when the limit exists. The annealed features are then
\begin{equation}
h^a_c(\gb) = \inf\{h \colon F^a(\gb,h) > 0\}, \quad \nu_a(\gb) = \lim_{h \searrow h_c^a(\gb)} \frac{\log F^a(\gb,h)}{\log (h - h_c^a(\gb))},
\end{equation}
when the limit exists (along the paper we may omit $\beta$ to lighten the notation, when there is no ambiguity). A simple application of Jensen's inequality leads to the following comparison:
\begin{equation}
F(\gb,h) \leq F^a(\gb,h) \text{ for all } h\in \R, \gb\ge0,
\end{equation}
and consequently
\begin{equation}
h_c^a(\gb) \leq h_c(\gb).
\end{equation}
If the annealed and quenched critical points or exponents differ at a given value of $\gb$, then disorder is said to be relevant for this value of $\beta$.

\subsubsection{The Harris criterion}

There have been a lot of studies on this problem in the past few years in the case when disorder is given by a sequence of i.i.d.~random variables with exponential moments (under this assumption the annealed model coincides with the homogeneous model after a suitable shift of $h$). All these works put the prediction known in the physics literature as the Harris criterion~\cite{H74} on a firm mathematical ground, which in this context states that disorder should be irrelevant if $\ga < 1/2$ (at least for small values of $\beta$) and relevant if $\ga > 1/2$. Several approaches have been used: direct estimates such as fractional moment and second moment estimates~\cite{AS06, AZ09, BL15, BCPSZ14, CadH13, DGLT09, GiT2006, GTL09}, martingale theory~\cite{L10}, variational techniques~\cite{CdH13}, and more recently chaos expansions of the partition functions~\cite{CSZ16+,CSZ16,CTT16+}. The limiting case $\ga = 1/2$ has been the subject of a lot of controversies and has been fully answered only recently~\cite{BL15}. Finally, the full criterion for relevance (in the sense of critical point shift) reads
\begin{equation}
\forall \gb > 0,\ h_c(\gb) > h_c^a(\gb) \iff \sum_{n \geq 1} L(n)^{- 2} n^{2 (\ga - 1)} = \infty,
\end{equation}
that corresponds to the intersection of two independent copies of $\tau$ being recurrent.

\subsubsection{Correlated disorder: state of the art}

The study of pinning models in correlated disorder is more recent, see~\cite{B13, B14, BL12, BP15, P13}. From a mathematical perspective, it is quite natural to try and understand how crucial the assumption of independence is for the basic properties of the polymer, and in particular for the validity of the relevance criterion. Also, in several instances, the sequence of inhomogeneities may present more or less strong correlations: let us mention for instance the sequence of nucleotides which play the role of the disorder sequence in DNA denaturation~\cite{JPS06}. The main idea is that the relevance criterion should be modified only if the correlations are strong enough. Note that with correlated disorder, even the annealed model may not be trivial. Mainly two types of correlated disorder have been considered until now: correlated Gaussian disorder~\cite{B13, BP15, P13} and random environments with large attractive regions of sub-exponential decay, also referred to as \emph{infinite disorder}~\cite{B14, BL12}.

\subsection{Scope of the paper}

The disorder sequence we consider is based on another renewal sequence $\htau$, independent of $\tau$, starting at the origin and whose law shall be denoted by $\hat\bP$. More specifically, we assume that if the Markov chain visits the origin at time $n$ then it is given a reward equal to one if $n \in \htau$, zero otherwise. We are therefore dealing with the following binary correlated disorder sequence:
\begin{equation}
\go_n = \hat\gd_n := \ind_{\{n \in \htau\}}, \quad n \in \bbN_0.
\end{equation}
From now on, the inter-arrival laws of $\tau$ and $\htau$ satisfy
\begin{equation}\label{ReturnTimeDistributions}
\begin{aligned}
K(n) &:= \bP(\tau_1 = n) \sim c_K \, n^{- (1 + \ga)},\\
\hK(n) &:= \hP(\htau_1 = n) = c_{\hK} \, n^{- (1 + \ha)}, \quad n \in \bbN
\end{aligned}
\end{equation}
with $\ga, \ha >0$, and
\begin{equation}
\mu := \bE(\tau_1),\qquad \hmu := \hE(\htau_1),
\end{equation}
which may be finite or infinite. Note that these definitions ensure aperiodicity for both renewal processes. In principle, the constants $c_K$ and $c_{\hK}$ may also be replaced by slowly varying functions, which would allow inclusion of the special case $\ga \in \{0,1\}$ in the discussion, but we refrain from doing so for the sake of simplicity. Also, we write an equality in the definition of $\hat K(n)$ to ensure log-convexity. This technical condition is actually only needed for proving Theorem~\ref{thm:smoothing} (see Lemmas~\ref{lem:mon.pro} and~\ref{lem:comp.tilt.shift}), which we actually believe to hold when the equality sign is replaced by the equivalent sign in the definition of $\hat K$ in~\eqref{ReturnTimeDistributions}.

The definitions of the basic thermodynamical quantities is the same as in the previous section, except that $\bbP$ and $\bbE$ are replaced by $\hat\bP$ and $\hat\bE$. The condition that $n \in \tau$ in the definition of the polymer measures above could be removed, leading to the \emph{free} versions. The versions with this condition are called the \emph{pinned} versions. It is a standard fact~\cite[Remark~1.2]{Gi2007} that this minor modification does not have any effect on the limiting free energies as defined in Propositions~\ref{pr:exist.ann.free.energy} and~\ref{pr:exist.que.free.energy}.

A first dichotomy arises:
\begin{itemize}
\item If $\ha < 1$, then the quantity in front of $\beta$ in the Hamiltonian is at most $|\htau \cap \{1,\ldots n\}|$, which is of order $n^{\ha}= o(n)$, and therefore disorder has no effect on the quenched free energy, which reduces to the homogeneous free energy. However the annealed model is non-trivial, so we include this case for completeness.

\item If $\ha > 1$, then (i) we may replace $\hP$ by its stationary version, denoted by $\hP_s$, under which the distribution of the increments $(\htau_{n + 1} - \htau_n)_{n \in \bbN_0}$ is the same as in $\hP$, whereas that of $\htau_0$ becomes $\{ \hP(\htau_1 > n) / \hmu\}_{n \in \bbN_0}$, see e.g.~\cite[Chapter~V, Corollary~3.6]{A2003} (again, this does not affect the free energy, see Propositions~\ref{pr:exist.ann.free.energy} and~\ref{pr:exist.que.free.energy}); (ii)~the correlation exponent of our environment is $\ha - 1$, since for $n > m$,
\begin{equation}\label{eq:speed_corr}
\begin{aligned}
\cov_{\hP_s}(\hat \gd_m, \hat \gd_n) &= \hat \bE_s(\hat \gd_m \hat \gd_n) - \hat \bE_s(\hat \gd_m) \hat \bE_s(\hat \gd_n)\\
&= \hP_s(m \in \htau) \Big(\hP_s(n \in \htau \mid m \in \htau) - \hP_s(n \in \htau) \Big)\\
&= \frac{1}{\hmu} \bigg(\hP(n - m \in \htau) - \frac{1}{\hmu} \bigg)\\
&\sim c (n - m)^{1 - \ha}, \qquad \text{ as } n - m \to \infty,
\end{aligned}
\end{equation}
for some positive constant $c$. The latter can be deduced from the Renewal Theorem and the following renewal convergence estimate~\cite[Lemma~4]{Fr82}
\begin{equation}\label{eq:ren.conv.est}
\hP(n \in \htau) - \frac{1}{\hmu} \sim \frac{c_{\hK}}{\ha (\ha - 1) \hmu^2} \frac{1}{n^{\ha - 1}}, \qquad n \to \infty.
\end{equation}
\end{itemize}

Although our choice of disorder may seem at first quite specific, it is motivated by the following:
\begin{itemize}
\item By tuning the value of the parameter exponent $\ha$, one finds a whole spectrum of correlation exponents ranging from non-summable correlations to sum\-mable correlations, according to whether the sum $\sum_{n \geq 0} \cov_{\hP_s}(\hat \gd_0, \hat \gd_n)$ is infinite or finite. According to~\eqref{eq:speed_corr}, correlations are summable when $\ha > 2$ and non-summable when $\ha < 2$.

\item Our disorder sequence is bounded, therefore the annealed free energy is always finite, in contrast to the case of Gaussian variables with non-summable correlations~\cite{B13}.

\item The probability of observing a long sequence of ones decays exponentially with length, which rules out the infinite disorder regime discussed in~\cite{B14}.

\item The renewal structure of the disorder sequence makes the study of the annealed model and decoupling inequalities more tractable.
\end{itemize}

\subsection{Summary of our results}

What we prove in this paper is the following: for the case $\ha > 2$ (summable correlations), disorder is irrelevant if $\ga < 1/2$ (both in the sense of critical points and critical exponents, at least for small values of $\beta$) and relevant if $\ga > 1/2$ (in the sense of critical exponents), which extends the Harris criterion for independent disorder. For the case $\ha \in (1,2)$ (non-summable correlations) all we are able to prove is that disorder is relevant when $\ga > 1/\ha$, a condition that we expect to be non-optimal. We discuss a list of predictions for disorder relevance in that case. Finally, in the case $\ha \in (0,1)$ disorder has no effect on the quenched free energy, but the annealed model exhibits a phase transition.

\subsection{Outline}

We present our results in Section~\ref{sec2}. Section~\ref{sec2.1} is dedicated to the annealed model and Section~\ref{sec2.2} to the quenched model. The proofs for the annealed model are in Section~\ref{sec3}, and for the quenched model in Section~\ref{sec4}. Results on renewal theory and homogeneous pinning are summarized in the appendix.


\section{Results}\label{sec2}

The intersection set of $\tau$ and $\htau$, which we denote by
\begin{equation}\label{eq:deftildetau}
\tilde\tau = \tau \cap \htau,
\end{equation}
will play a fundamental role in the remainder of the paper. Note that it is itself a renewal starting at $\tilde\tau_0 = 0$. We denote its law by $\tilde\bP$ and write
\begin{equation}\label{eq:deftildedelta}
\tilde\gd_n := \ind_{\{n \in \tau \cap \htau\}} = \ind_{\{n \in \tau\}} \ind_{\{n \in \htau\}} = \gd_n \hat \gd_n, \qquad n \in \bbN_0.
\end{equation}

\subsection{Results on the annealed model} \label{sec2.1}

We begin with the existence of the annealed free energy.

\begin{prop}\label{pr:exist.ann.free.energy}
For all $\gb \geq 0$ and $h \in \bbR$, the annealed free energy
\begin{equation}
F^a(\gb,h) = \lim_{n \to \infty} \frac{1}{n} \log Z^a_{n,\gb,h}
\end{equation}
exists and it is finite and non-negative. The result still holds, without changing the value of the free energy, when $\hmu < \infty$ and $\hP$ is replaced by $\hP_s$.
\end{prop}

The following basic properties of the annealed free energy are standard: the function $(\gb,h) \mapsto F^a(\gb,h)$ is convex, continuous, and non-decreasing in both variables.

\subsubsection{An auxiliary function: the number of intersection points}

Before stating further results, we need to introduce an auxiliary function which will help us to characterize the annealed critical point. For $h \leq 0$, denote by $\bP_h$ the probability of the transient renewal process with $\tau_0 = 0$ and inter-arrival law
\begin{equation}\label{hTransRenewal}
K_h(n) = e^h K(n), \quad n \in \bbN, \qquad K_h(\infty) = 1 - e^h.
\end{equation}
We denote the corresponding expectation by $\bE_h$. The expected number of points in the renewal set $\tilde\tau$ (including $0$) under the law $\bP_{h} \times \hP$ is denoted by
\begin{equation}
\cI(h) := \bE_h \hE(|\tilde\tau|) \in [1,\infty].
\end{equation}
Note that
\begin{equation}\label{eq:cI}
\begin{aligned}
\cI(h) &= \sum_{n \in \bbN_0} \bP_h(n \in \tau) \hP(n \in \htau) = \sum_{n,k \in \bbN_0} e^{h k} \bP(\tau_k = n) \hP(n \in \htau)\\
&= \sum_{k \in \bbN_0} e^{h k} \bP \times \hP(\tau_k \in \htau).
\end{aligned}
\end{equation}
The function $\cI$ is finite and infinitely differentiable in $(- \infty, 0)$. It is also continuous in $(- \infty, 0]$, increasing, and strictly convex. Its range is $[1,\cI(0)]$ with $\cI(0) = \bE \hE(|\tau \cap \htau|)$, which may be finite or infinite. It follows from Proposition~\ref{meanReturns} that
\begin{equation}\label{eq:rel_p_I}
p(h) := \bP_h \times \hP(\tilde\tau_1 < \infty) = 1 - \cI(h)^{-1}.
\end{equation}

Our next result provides an expression for the annealed critical curve involving the function $\cI$.

\begin{prop}\label{pr:ann.crit.curve}
Let $\gb_0 = - \log~(0)$. The annealed critical curve is
\begin{equation}\label{anncc}
h^a_c(\gb) =
\begin{cases}
\cI^{-1} \big(\frac{1}{1 - e^{-\gb}} \big) & \text{ if } \gb > \gb_0, \\
0 & \text{ if } 0 \leq \gb \leq \gb_0.
\end{cases}
\end{equation}
\end{prop}



The first term $- \gb / \hmu$ simply accounts for the fact that our disorder sequence is not centered, and that by the Renewal Theorem, $\lim_{n \to \infty} \hE(\hat \gd_n) = 1 / \hmu$. Note that by Jensen's inequality, $h_c^a(\gb) \leq - \gb / \hmu$, and this already gives that $c_a \geq 0$ in Proposition~\ref{pr:ann.cc.scaling}. If $\ha > 1 + \ga \wedge 1$, then $\gann = 2$, as in the i.i.d.~case, but if $\ha < 1 + \ga \wedge 1$, there is an anomalous scaling of the annealed critical curve. Moreover, if $\ha <1$ then $\hmu = \infty$, so the term $\gb/\hmu$ disappears and $\gann >1$ gives the first-order term.


\begin{remark}\label{rmk:signbeta0}
From~\eqref{eq:rel_p_I} we have that $\gb_0 = - \log(1 - \{\bE \hE(|\tilde\tau|)\}^{-1})$ is non-negative. Therefore, using Proposition~\ref{RMassFunctionProp}, we see that
\begin{equation}\label{eq:logp0}
\gb_0 \,
\begin{cases}
>0 & \text{ if } \ga + \ha < 1,\\
=0 & \text{ if } \ga + \ha > 1.
\end{cases}
\end{equation}
\end{remark}

\begin{prop}\label{pr:ann.cc.scaling}
Suppose $\ga + \ha > 1$ (then $\gb_0 = 0$). There exists $c_a > 0$ such that
\begin{equation}
h_c^a(\gb) = - \frac{\gb}{\hmu} - c_a \gb^{\gann}[1 + o(1)], \quad \text{ as }\gb \searrow 0,
\end{equation}
where
\begin{equation}\label{eq:defgann}
\gann =
\begin{cases}
1 + \left[\frac{\ha -1}{\ga \wedge 1} \wedge 1 \right] & \text{if } \ha > 1 \text{and } \ha \neq 1 + \ga \wedge 1\\[0.3em]
\frac{\ga \wedge 1}{\ha -1 + \ga \wedge 1} & \text{if } \ha < 1.
\end{cases}
\end{equation}
If $\ha = 1 + \ga \wedge 1$, we get instead
\begin{equation}
h_c^a(\gb) = - \frac{\gb}{\hmu} - c_a \gb^2 |{\log \gb}|[1 + o(1)].
\end{equation}
\end{prop}

\begin{prop}\label{pr:ann.cc.scaling2}
Suppose $\ga + \ha < 1$ (then $\gb_0 > 0$). As $\gb \searrow \gb_0$, there is a constant $c \in (- \infty,0)$ such that
\begin{equation}
h_c^a(\gb) \sim c\frac{(\gb - \gb_0)^{\gann}}{1+|{\log(\gb-\gb_0)}|\ind_{\{1-\hat \ga=2\ga\}}}, \quad \text{where} \quad
\gann = 1 \vee \frac{\ga}{1 - \ga - \ha}.
\end{equation}
\end{prop}

\section{Proof of the annealed results}\label{sec3}

The main idea is that the annealed model can be viewed as a homogeneous pinning model for the intersection renewal $\tilde\tau = \tau \cap \htau$ after the law of $\tau$ has been tilted.

\subsection{Existence of the free energy}\label{sec3.1}

\begin{proof}[Proof of Proposition~\ref{pr:exist.ann.free.energy}]
We use standard techniques, see the proof of Lem\-ma~3.5 in~\cite{Gi2011}. Let us first introduce the \emph{fully-pinned} annealed partition function
\begin{equation}\label{eq:fullpinned}
Z^{a, c}_{n,\gb,h} := \hat \bE \bE \left(e^{\sum_{k=1}^n (h + \gb \hat \gd_k) \gd_k} \gd_n \hat \gd_n \right).
\end{equation}
We shall write $Z^{a, c}_{n,\gb,h}(A)$ when the expectation above is restricted to the event $A$ and use the same convention for other versions of the partition function appearing in the proof. By super-additivity, the sequence $\{n^{-1} \log Z_{n,h,\gb}^{a,c}\}_{n \in \bbN}$ converges to the limit
\begin{equation}\label{eq:facsup}
F^{a, c}(\gb,h) := \sup_{n \in \bbN} \{n^{-1} \log Z^{a,c}_{n,\gb,h}\}.
\end{equation}
From the bounds
\begin{equation}
e^{|h + \gb| n} \geq Z^{a,c}_{n,\gb,h} \geq Z^{a,c}_{n,\gb,h}(\tau_1 = \htau_1 = n) = e^{h + \gb} K(n) \hK(n),
\end{equation}
we get that $F^{a,c}(\gb,h) \in [0,\infty)$. Let us now prove that $F^a(\gb,h)$ exists and that $F^a(\gb,h) = F^{a,c}(\gb,h)$. Let $R = \sup \{k \leq n \colon k\in \tilde\tau\}$ be the last point in $\tilde\tau$ before $n$. Then,
\begin{equation}
\begin{aligned}
\!\!\!Z^{a}_{n,\gb,h} &= \hat \bE \bE \!\left(e^{\sum_{k=1}^n (h + \gb \hat \gd_k) \gd_k} \gd_n \right)\\
&= \sum_{r=0}^n \hat \bE \bE \!\left(e^{\sum_{k=1}^r (h + \gb \hat \gd_k) \gd_k} e^{\sum_{k=r+1}^n (h + \gb \hat \gd_k) \gd_k} \ind_{\{R = r\}} \gd_n \right)\\
&= \sum_{r=0}^n \hat \bE \bE \!\left(e^{\sum_{k=1}^r (h + \gb \hat \gd_k) \gd_k} \gd_{r} \hat \gd_r \right) \!\hat \bE \bE\! \left(e^{\sum_{k=r+1}^n (h + \gb \hat \gd_k) \gd_k} \ind_{\{\tilde{\tau} \cap [r+1,n] = \emptyset\}} \gd_n \,\middle|\, r \Mk\in\Mk \tilde\tau \right)\!\!\\
&= \sum_{r=0}^n Z_{r,\gb,h}^{a,c} \hat \bE \bE \!\left(e^{h \sum_{k=r+1}^n \gd_k}\ind_{\{\tilde{\tau} \cap [r+1,n] = \emptyset\}} \gd_n \,\middle|\, r \in \tilde\tau \right).
\end{aligned}
\end{equation}
From~\eqref{eq:facsup} we know that $Z_{r,\gb,h}^{a,c} \leq e^{rF^{a,c}(\gb,h)}$. Moreover,
\begin{multline}
\hat \bE \bE \left(e^{h \sum_{k=r+1}^n \gd_k} \ind_{\{\tau \cap \htau \cap [r+1,n] = \emptyset\}} \gd_n \,\middle|\, r \in \tilde\tau \right)\\
\begin{aligned}
&\leq \hat \bE \bE \left(e^{h \sum_{k=r+1}^n \gd_k} \gd_n \,\middle|\, r \in \tilde\tau \right)\\
&= Z^{a,c}_{n-r,0,h} \frac{1}{\hP(n - r \in \htau)} \leq e^{(n - r) F^{a,c}(0,h)} \frac{1}{\hP(n - r \in \htau)}\\
&\leq e^{(n - r)F^{a,c}(\gb,h)} \frac{1}{\hP(n - r \in \htau)}.
\end{aligned}
\end{multline}
Thus,
\begin{equation}\label{eq:zna_ub}
Z^a_{n,\gb,h} \leq \sum_{r=0}^n e^{r F^{a,c}(\gb,h)} e^{(n - r) F^{a,c}(\gb,h)} \frac{1}{\hP(n - r \in \htau)} = e^{n F^{a,c}(\gb,h)} \sum_{r=0}^n \frac{1}{\hP(r \in \htau)}.
\end{equation}
By Proposition~\ref{RMassFunctionProp}, the last sum increases polynomially in $n$. Combining this inequality with $Z^{a,c}_{n,\gb,h} \leq Z^a_{n,\gb,h}$, we get that the free energy $F^a(\gb,h)$ exists and equals $F^{a,c}(\gb,h)$.

We now prove the second part of the result, namely that the limit for $\hP_s$ is the same as for $\hP$. Suppose that $\hmu < \infty$, and define $Z_n^{a,s} = \bE \hE_s \left(\exp \left\{ \sum_{k=1}^n (h + \gb \hat \gd_k) \gd_k \right\} \gd_n \right)$ (we temporarily remove $\gb$ and $h$, for conciseness). By restricting the expectation to the event $\{0 \in \htau\}$, we get on the one hand $Z_n^{a,s} \geq Z_n^a / \hmu$. On the other hand, by decomposing on the value of $\tilde\tau_1$, we obtain
\begin{equation}\label{eq:znas_ub}
\begin{aligned}
Z_n^{a,s} &= Z_n^{a,s}(\tilde\tau_1 > n) + \sum_{k=1}^n Z_n^{a,s}(\tilde\tau_1 = k)\\
&\leq e^{n (F(0,h) + o(1))} + \sum_{k=1}^n e^{k F(0,h)} Z^a_{n - k}.
\end{aligned}
\end{equation}
To go from the first to second line, we used the Markov property at $\tilde\tau_1$ and the fact that $\gd_k \hat \gd_k = 0$ for all $k < \tilde\tau_1$. Combining~\eqref{eq:znas_ub}, \eqref{eq:zna_ub}, and using that $F(0,h) \leq F^a(\gb,h)$, we get that
\begin{equation}
\limsup_{n \to \infty} \frac{1}{n} \log Z_n^{a,s} \leq F^a(\gb,h),
\end{equation}
hence the result.
\end{proof}

\subsection{Annealed critical curve}\label{sec3.2}

This subsection is organized as follows: we start with Lemma~\ref{lem:acc} below, which we use to prove Proposition~\ref{pr:ann.crit.curve}. From Lemma~\ref{lem:asympt.0} we get Proposition~\ref{pr:ann.cc.scaling2} and Lemma~\ref{lem:asympt.I}, which in turn yields Proposition~\ref{pr:ann.cc.scaling}.

\begin{lemm}\label{lem:acc}
Let $h \leq 0$. Then, $F^a(\gb,h) = 0$ if and only if $\gb \leq - \log p(h)$.
\end{lemm}

\begin{proof}[Proof of Lemma~\ref{lem:acc}]
Let $h \leq 0$. We know from the proof of Proposition~\ref{pr:exist.ann.free.energy} that $F^a(\gb,h)$ is the limiting free energy of the fully-pinned partition function in~\eqref{eq:fullpinned}, which we may rewrite, using~\eqref{hTransRenewal} and~\eqref{eq:deftildedelta}, as
\begin{equation}\label{eq:equ.znac}
Z^{a,c}_{n,\gb,h} = \hE \bE_h \left(e^{\gb \sum_{k=1}^n \tilde \gd_k} \tilde \gd_n \right),
\end{equation}
which is the partition function of the homogenous pinning model with reward $\gb$ for the renewal $\tilde\tau$ under the law $\bP_h \times \hP$. It now follows from a standard fact about homogeneous pinning models (see~\cite[Section~1.2.2, Equation~(1.26)]{Gi2007}) that its critical point (as $\gb$ varies and $h$ is fixed) is at
\begin{equation}
\gb = - \log \bP_h \hP(\tilde \tau_1 < \infty) = - \log p(h),
\end{equation}
see~\eqref{eq:rel_p_I}.
\end{proof}

\begin{proof}[Proof of Proposition~\ref{pr:ann.crit.curve}]
We distinguish two cases. Suppose first that $\gb \leq \gb_0 = - \log~(0)$. On the one hand, if $h \leq 0$, then $F^a(\gb,h) \leq F^a(\gb,0) = 0$, by Lemma~\ref{lem:acc}. On the other hand, if $h > 0$ then $F^a(\gb,h) \geq F^a(0,h) > 0$, by~\cite[Theorem~2.7]{Gi2007}, since $\tau$ is recurrent. Thus, $h_c^a(\gb) = 0$ in this case. Suppose now that $\gb > \gb_0$. If $h \geq 0$, $F^a(\gb,h) > 0$ for the same reason as above, so we restrict to $h < 0$. Then, by Lemma~\ref{lem:acc}, $F^a(\gb,h) = 0$ if and only if $\gb \leq - \log p(h)$, that is $\cI(h) \geq (1 - e^{- \gb})^{-1}$ (recall~\eqref{eq:rel_p_I} and the lines above). Since $\gb < \gb_0$, we have $\cI(0)=(1 - e^{- \gb_0})^{-1}> (1 - e^{- \gb})^{-1} > 1$, that is, $(1 - e^{- \gb})^{-1}$ is in the range of $\cI$, so we get
\begin{equation}
\gb \leq - \log p(h) \quad \text{if and only if } \quad h \geq \cI^{-1} \bigg(\frac{1}{1 - e^{- \gb}} \bigg),
\end{equation}
hence the result.
\end{proof}

\begin{proof}[Proof of Proposition~\ref{pr:ann.cc.scaling}]
Let us first consider the case $\ha > 1$ and $\ha \neq 1 + \ga \wedge 1$, and write
\begin{equation}\label{eq:asymphca}
h_c^a(\gb) = - \frac{\gb}{\hmu} (1 + \gep_\gb) \qquad \text{with } \lim_{\gb\downarrow0} \gep_\gb = 0.
\end{equation}
Note that $\ha > 1$ implies $\gb_0 = 0$, by~\eqref{eq:logp0}. Therefore, from Proposition~\ref{pr:ann.crit.curve}, we get on the one hand
\begin{equation}\label{eq:as.hca1}
\cI(h_c^a(\gb)) = \frac{1}{1 - e^{- \gb}} = \frac{1}{\gb} \bigg(1 + \frac{1}{2} \gb + o(\gb) \bigg), \qquad \gb \downarrow 0,
\end{equation}
and on the other hand, from Lemma~\ref{lem:asympt.I} below and~\eqref{eq:asymphca}, as $\gb \downarrow 0$,
\begin{equation}\label{eq:as.hca2}
\begin{aligned}
\cI(h_c^a(\gb)) &= \frac{1}{\hmu} \frac{1}{1 - e^{h_c^a(\gb)}} + c \, (- h_c^a(\gb))^{\gann - 2} [1 + o(1)]\\
&= \frac{1}{\gb} \bigg(1 - \gep_\gb [1 + o(1)] + \frac{1}{2 \hmu} \gb [1 + o(1)] + c \, \gb^{\gann - 1} [1 + o(1)] \bigg),
\end{aligned}
\end{equation}
where $c$ is a positive constant that may change from line to line and $\gann$ is defined as in~\eqref{eq:defgann}. The result follows by identifying the right-hand sides in~\eqref{eq:as.hca1} and~\eqref{eq:as.hca2}. Indeed we get $\gep_\gb \sim c \gb^{\gann - 1}$ when $\gann<2$ and $\gep_\gb \sim (c-\tfrac12 (1-1/\hmu)) \gb$ when $\gann =2$. Note that in the latter case the constant $c-\tfrac12 (1-1/\hmu)$ is indeed positive by Lemma~\ref{lem:asympt.I}. If $\ha = 1 + \ga \wedge 1$, the same method leads us to $\gep_\gb \sim c \gb |\log \gb|$, which proves our claim. Finally, the case $\ha < 1$ is easier. Indeed, \eqref{eq:as.hca1} gives $\cI(h_c^a(\gb)) \sim 1/\gb$ and the result follows from Lemma~\ref{lem:asympt.I}.
\end{proof}

\begin{proof}[Proof of Proposition~\ref{pr:ann.cc.scaling2}]
Since $\gb_0 > 0$ we have by Proposition~\ref{pr:ann.crit.curve}
\begin{equation}\label{eq:ann.cc2a}
\cI(h_c^a(\gb_0)) - \cI(h_c^a(\gb_0 + \gep)) \sim \frac{e^{-\gb_0}}{(1-e^{-\gb_0})^2} \gep, \qquad \gep\to0.
\end{equation} Moreover, for some positive constant $c>0$,
\begin{equation}\label{eq:ann.cc2b}
\cI(0) - \cI(h) = \sum_{k\in\bbN_0} (1- e^{h k}) \bE \hP(\tau_k \in \htau) \sim c |h|^{1\wedge \frac{1-\ha-\ga}{\ga}}(1+|{\log |h|}|\ind_{\{1-\hat \ga=2\ga\}})
\end{equation}
as $h\to 0^-$. The equivalence above follows from Lemma~\ref{seriesAsymptotics}. We get the final result by noting that $h_c^a(\gb_0)=0$ and combining~\eqref{eq:ann.cc2a} and~\eqref{eq:ann.cc2b}.
\end{proof}

\begin{lemm}\label{lem:asympt.0}
\begin{enumerate}\romanenumi
\item \label{3.2.1} If $\ha > 1$ then, as $k \to \infty$,
\begin{equation}
\bE \hP (\tau_k \in \htau) - \frac{1}{\hmu} \sim \frac{c_{\hK}}{\hmu^2 \ha (\ha - 1)} \,
\begin{cases}
\mu^{1 - \ha} k^{1 - \ha} & \text{if } \ga > 1,\\[2pt]
\bE(X_\ga^{1 - \ha}) k^{\frac{1 - \ha}{\ga}} & \text{if } \ga \in (0,1),
\end{cases}
\end{equation}
where $X_\ga$ is an $\ga$-stable random variable totally skewed to the right, with scale parameter $\gs > 0$ depending on the distribution of $\tau_1$ and shift parameter $0$ (see relation~\eqref{eq:stable.cf} in the appendix for a reminder of these terms).

\item \label{3.2.2} If $\ha \in (0,1)$,
\begin{equation}
\bE \hP (\tau_k \in \htau) \sim \frac{C_{\ha}}{c_{\hK}} \,
\begin{cases}
\mu^{\ha - 1} k^{\ha - 1} & \text{if } \ga > 1,\\[2pt]
\bE(X_\ga^{\ha - 1}) k^{\frac{\ha - 1}{\ga}} & \text{if } \ga \in (0,1),
\end{cases}
\end{equation}
where $C_{\ha}$ is as in Proposition~\ref{RMassFunctionProp}.
\end{enumerate}
\end{lemm}

\begin{proofc}[Proof of Lemma~\ref{lem:asympt.0}]

\begin{step}{\eqref{3.2.1}}
From the renewal convergence estimates in~\cite[Lemma~4]{Fr82}, for $\ha > 1$, we get
\begin{equation}
\hP(n \in \htau) - \frac{1}{\hmu} \sim \frac{c_{\hK}}{\hmu^2 \ha (\ha - 1)} n^{1 - \ha}, \quad n \to \infty,
\end{equation}
so we have $\bP$-a.s.,
\begin{equation}
\hP(\tau_k \in \htau) - \frac{1}{\hmu} \sim \frac{c_{\hK}}{\hmu^2 \ha (\ha - 1)} \tau_k^{1 - \ha}, \quad k \to \infty.
\end{equation}
Since $\tau_k \geq k$, we may take the expectation in the line above and write
\begin{equation}
\bE \hP(\tau_k \in \htau) - \frac{1}{\hmu} \sim \frac{c_{\hK}}{\hmu^2 \ha (\ha - 1)} \bE(\tau_k^{1 - \ha}), \quad k \to \infty,
\end{equation}
and we may conclude the proof with Lemma~\ref{lem:conv.power.tauk}.
\end{step}

\begin{laststep}{\eqref{3.2.2}}
From Proposition~\ref{RMassFunctionProp}, we have $\bP$-a.s, if $\ha \in (0,1)$,
\begin{equation}
\hP(\tau_k \in \htau) \sim \frac{C_{\ha}}{c_{\hK}} \tau_k^{\ha - 1},\quad k \to \infty,
\end{equation}
and with the same argument as in~\eqref{3.2.1},
\begin{equation}
\bE \hP(\tau_k \in \htau) \sim \frac{C_{\ha}}{c_{\hK}} \bE(\tau_k^{\ha - 1}), \quad k \to \infty.
\end{equation}
We may conclude thanks to Lemma~\ref{lem:conv.power.tauk}.
\end{laststep}
\let\qed\relax
\end{proofc}

\begin{lemm}\label{lem:asympt.I}
Suppose that $\ga + \ha > 1$. As $h \uparrow 0^-$,
\begin{equation}
\cI(h) - \frac{1}{\hmu} \frac{1}{1 - e^h} \sim 
\begin{cases}
c|h|^{\frac{1-\ha}{\ga\wedge 1} - 1} & \text{if } \ha \in (1-\ga\wedge 1, 1),\\
c|h|^{\frac{\ha - 1}{\ga\wedge 1} - 1} & \text{if } \ha \in (1, \ga \wedge 1+1),\\
c|{\log |h|}| & \text{if } \ha = \ga \wedge 1+1,\\
c & \text{if } \ha > \ga \wedge 1+1,
\end{cases}
\end{equation}
where $c$ is a positive constant (note that in the first case $\hmu = \infty$ and the left-hand side is simply $\cI(h)$). Moreover, in the case $ \ha > \ga \wedge 1+1$, the constant $c$ satisfies $c>(1-\hat \mu^{-1})/2$.
\end{lemm}

\begin{proof}[Proof of Lemma~\ref{lem:asympt.I}]
Recall~\eqref{eq:cI}. For $h < 0$, we may write
\begin{equation}
\cI(h) - \frac{1}{\hmu} \frac{1}{1 - e^h} = \sum_{k\in\bbN_0} e^{h k} \bE \bigg[\hP(\tau_k \in \htau) - \frac{1}{\hmu} \bigg].
\end{equation}
Then, in the first three cases, the result follows from Lemma~\ref{lem:asympt.0} and the standard Tauberian arguments recalled in Lemma~\ref{seriesAsymptotics}$\MK$\eqref{A.8.3}. In the fourth case, we have
\[
c=\sum_{k\in\bbN_0} \bE \bigg[\hP(\tau_k \in \htau) - \frac{1}{\hmu} \bigg].
\]
Note that the set $\sigma:=\{k\in \bbN_0: \tau_k\in \htau\}$ defines a renewal process. Call $\gs_1$ its first positive point. This has mean $\hat \mu$ because $\bP\times \hP(k\in \sigma)\to 1/\hat \mu$ as $k\to\infty$. Let us call $v\in(0, \infty]$ its variance. We now use Problem 19 in Chapter~XIII of~\cite{Feller68}, which needs a fix: we let the reader check that in equations~(12.1) and~(12.2) therein, $u_0$ should be removed and the sums should start at $n=0$. According to this, we have
\[
c=\sum_{k\in\bbN_0} \bE \bigg[\hP(k \in \gs) - \frac{1}{\hmu} \bigg]=\frac{v^2+\hat \mu^2-\hat \mu}{2\hat \mu^2},
\]
so that $c-(1-\hat \mu^{-1})/2=v^2/(2\hat\mu^2)>0$.
\end{proof}

\section*{Appendix. Results on renewals and homogeneous pinning} \label{Facts}
\setcounter{section}{1}
\setcounter{equation}{0}
\setcounter{cdrthm}{0}

We collect here a few results on renewal processes and the homogeneous pinning model.

\begin{prop}[{See~\cite[Proposition~2.4, Chapter~I]{A2003}}] \label{meanReturns}
If $\tau$ is a (possibly transient) renewal with $\tau_0 = 0$, then $\bE(|\tau|) = \bP(\tau_1 = \infty)^{-1}$.
\end{prop}

\begin{prop}[{See~\cite[Theorem~A.4, Appendix~A.5.2]{Gi2007}}]
If $\tau$ is a transient renewal that satisfies~\eqref{eq:defK_regvar} then
\begin{equation}
\bP(n \in \tau) \sim \frac{\bP(\tau_1 = n)}{\bP(\tau_1 = \infty)^2} \qquad \text{as } n \to \infty.
\end{equation}
\end{prop}

\begin{prop}\label{RMassFunctionProp}
If $\tau$ is a recurrent renewal that satisfies~\eqref{eq:defK_regvar} then as $n \to \infty$,
\begin{equation}\label{RMassFunction}
\bP(n \in \tau) \sim
\begin{cases}
\frac{K(n)}{(\bar K(n))^2} & \text{ if } \ga = 0,\\[4pt]
\frac{C_\ga}{L(n) n^{1 - \ga}} & \text{ if } \ga \in (0,1),\\[4pt]
\left\{ \sum_{j=0}^n \bar K(j) \right\}^{-1} & \text{ if } \ga \geq 1,
\end{cases}
\end{equation}
with $\bar K(j) = \sum_{r=j+1}^\infty K(r)$ for each $j\in\bbN_0$ and $C_\ga = \ga \sin(\pi \ga)/\pi$.
\end{prop}

When $\ga = 0$, the sequence $(\bar K(n))_{n\in\bbN}$ is slowly varying (\cite[Propos-\linebreak ition~1.5.9$\MK$(b)]{BGT89}), while when $\ga = 1$, the sequence $(\sum_{j=0}^ n \bar K(j))_{n\in\bbN}$ is slowly varying (\cite[Prop\-os\-itions~1.5.10 and~1.5.9$\MK$(a)]{BGT89}).

The proof of the statement for $\ga = 0$ is in~\cite[Theorem~1.1]{Na2012}, for $\ga \in (0,1)$ in~\cite[Theorem~1.1]{GL63}, for $\ga = 1$ in~\cite[Theorem~8.7.5]{BGT89}, while for $\ga > 1$ it is the Renewal Theorem.

\begin{prop}\label{pr:transient.inter}
Let $\tau$ and $\tau'$ be two independent copies of a recurrent renewal process that satisfies~\eqref{eq:defK_regvar}. If $\ga < \frac{1}{2}$, or $\ga = \frac{1}{2}$ and $\sum_{n\geq1} n^{-1} L(n)^{-2} < \infty$, then $\tau \cap \tau'$ is transient.
\end{prop}

\begin{proof}
Use that $\bE^{\otimes 2}(|\tau \cap \tau'|) = \sum_{n\geq0} \bP(n \in \tau)^2$, in combination with Prop\-os\-ition~\ref{RMassFunctionProp}.
\end{proof}

Before stating the next lemma, we recall that an $\ga$-stable random variable $X$ has three parameters $\kappa\in[-1,1]$, $\sigma\ge 0$ and $m\in \bbR$ (skewness, scale and shift, respectively) appearing in its characteristic function:
\begin{equation} \label{eq:stable.cf}
\bE(e^{i\theta X}) =
\begin{cases}
\exp(-\sigma |\theta| (1 + i\kappa (2/\pi) \sign(\theta) \log|\theta|) + im\theta) & \text{if }\ga = 1,\\[4pt]
\exp(-\sigma^\ga |\theta|^\ga (1 - i\kappa \sign(\theta) \tan(\pi \ga /2)) + im\theta) & \text{if } \ga\in(0, 1)\cup(1,2),
\end{cases}
\end{equation} see Definition~1.1.6 in~\cite{ST94}. If $\kappa =1$ (resp.\ $-1$), $X$ is said to be totally skewed to the right (resp.\ left).

\begin{lemm}\label{lem:conv.power.tauk}
If $r > 0$, then
\begin{equation}
\bE(\tau_k^{- r}) \sim
\begin{cases}
(\mu k)^{- r} & \text{if } \ga > 1,\\[4pt]
\bE(X_\ga^{- r}) k^{- r/\ga} & \text{if } \ga \in (0,1),
\end{cases}
\qquad k \to \infty,
\end{equation}
where $X_\ga$ is an $\ga$-stable random variable totally skewed to the right, with scale parameter $\gs > 0$ depending on the distribution of $\tau_1$ and shift parameter $0$.
\end{lemm}

\begin{proof}
If $\ga > 1$, the result follows by bounded convergence since, by the Renewal Theorem, $(\tau_k/k)^{- r}$ converges $\bP$-a.s.~to $\mu^{- r}$ and is bounded from above by $1$. If $\ga \in (0,1)$, we use that $\tau_k/k^{1/\ga}$ converges to an $\ga$-stable random variable $X_\ga$. The only complication is that $(\tau_k/k^{1/\ga})^{- r}$ is not bounded, but the result still holds by uniform integrability, namely, by Exercise~3.2.5 in~\cite{Du10}. It is enough to show that for some $\gga > r$ we have
\begin{equation}\label{eq:UItau}
\sup_{k\geq1} \bE(\{\tau_k/k^{1/\ga}\}^{- \gga}) < \infty.
\end{equation}
To show this, first note that
\begin{equation}\label{ExpectationTail}
\bE(\{\tau_k/k^{1/\ga}\}^{- \gga}) = \int_0^\infty \bP(\{\tau_k/k^{1/\ga}\}^{- \gga} > t) \, \dd t = \int_0^\infty \bP(\tau_k < k^{1/\ga} t^{- 1/\gga}) \, \dd t.
\end{equation}
With the use of Chernoff's bound, the probability inside the integral is bounded as
\begin{equation}\label{LDPBound}
\bP(\tau_k/k < k^{(1/\ga) - 1} t^{- 1/\gga}) \leq \exp \Big\{ - k \sup_{\gl\leq0} \{\gl k^{(1/\ga) - 1} t^{- 1/\gga} - \log M(\gl)\} \Big\},
\end{equation}
where $M(\gl) := \bE(e^{\gl \tau_1})$. A standard Tauberian argument~\cite[(5.22) of Chapter~XIII]{Feller71} shows that there exists a constant $C > 0$ such that $M(\gl) \leq \exp({- C |\gl|^\ga})$ for all $\gl \leq 0$. This implies, for $x > 0$, the following bound
\begin{equation}
\sup_{\gl\leq0} \{\gl x - \log M(\gl)\} \geq \sup_{\gl\leq0} \{\gl x + C |\gl|^\ga\} = C_1 x^{- \ga/(1 - \ga)},
\end{equation}
where $C_1 := (1 - \ga) C^{(1 - \ga)^{-1}} \ga^{\ga (1 - \ga)^{-1}} > 0$. Consequently, the probability in~\eqref{LDPBound} is bounded from above by $\exp\{- C_1 t^\frac{\ga}{(1 - \ga) \gga}\}$, which completes the proof.
\end{proof}

\begin{lemm}\label{seriesAsymptotics}
Let $(r_k)_{k\ge1}$ be a sequence of positive numbers so that $r_k\sim k^{-\gl}$ as $k\to\infty$ for some $\gl>0$. Then
\begin{equation} \tag{i} \label{A.8.1}
\sum_{k=1}^\infty (1-e^{-xk})r _k\sim
\begin{cases} x A_1& \text{ if } \gl>2, \\
x |{\log x}| & \text{ if } \gl=2, \\
x^{\gl-1}A_2(\gl) & \text{ if } \gl\in(1, 2) \\
\end{cases}
\end{equation}
as $x\to 0^+$, where $A_1=\sum_{k=1}^\infty k r_k$ and $A_2(\gl)=\int_0^\infty (1-e^{-t}) t^{-\gl}\, \dd t$ are positive and finite.
\begin{equation} \tag{ii} \label{A.8.2}
\sum_{k=1}^\infty (e^{-xk}-1+xk) r_k\sim
\begin{cases} x^2 A_3& \text{ if } \gl>3, \\
x^2|{\log x}|/2 & \text{ if } \gl=3, \\
x^{\gl-1}A_4(\gl) & \text{ if } \gl\in(2, 3) \\
\end{cases}
\end{equation}
as $x\to 0^+$, where $A_3=(1/2)\sum_{k=1}^\infty k^2 r_k$ and $A_4(\gl)=\int_0^\infty (e^{-t}-1+t) t^{-\gl}\, \dd t$ are positive and finite.
\begin{equation} \tag{iii} \label{A.8.3}
\sum_{k=1}^\infty e^{-xk} r_k \sim
\begin{cases}
|{\log x}| & \text{ if } \gl = 1,\\[2pt] A_5(\gl) x^{\gl-1}& \text{ if } \gl \in (0,1)
\end{cases}
\end{equation}
as $x\to 0^+$, where $A_5(\gl)=\int_0^\infty e^{-t}t^{-\gl}\, \dd t$ is positive and finite.
\end{lemm}
%\let\theequation\oldtheequation

\begin{proof}
References for $\gl\in (1,2]$ in~\eqref{A.8.1} and $\gl\in(0,1]$ in~\eqref{A.8.3} are Corollary~8.1.7 and Theorem~1.7.1 in~\cite{BGT89}, respectively. If $\gl>2$ in~\eqref{A.8.1}, we have
\[
\frac{1}{h}\sum_{k=1}^\infty (1-e^{-hk}) r_k=\sum_{k=1}^\infty \frac{1-e^{-hk}}{hk} k r_k
\]
and to the last sum we apply the Dominated Convergence Theorem since $(1-e^{-x})/x$ is positive and bounded from above by 1 while $\sum_{k=1}^\infty k r_k<\infty$. Case~\eqref{A.8.2} follows from similar arguments and is left to the reader.
\end{proof}

In the following two lemmas we denote by $\bP_\gt$ the renewal such that $\bP_\gt(\tau_1 = n) =\linebreak e^{\gt - F(\gt) n} K(n)$, where $\gt\le 0$. Note that $\bP_0$ coincides with $\bP$. We also define $\phi_\gt(t) = \bE_\gt(e^{i t \tau_1})$, for $t\in\bbR$.

\begin{thebibliography}{MMMMMM}
\bibitem[AB18]{A17_1} Alexander, Kenneth S.; Berger, Quentin Pinning of a renewal on a quenched renewal, Electron. J. Probab., Volume 23 (2018), 6, 48 pages \href{https://ahl.centre-mersenne.org/item/AHL_2019__2__281_0/#A17_1}{Publisher reference}.
\bibitem[AS06]{AS06} Alexander, Kenneth S.; Sidoravicius, Vladas Pinning of polymers and interfaces by random potentials, Ann. Appl. Probab., Volume 16 (2006) no. 2, pp. 636-669 \href{https://ahl.centre-mersenne.org/item/AHL_2019__2__281_0/#AS06}{Publisher reference}.
\bibitem[Asm03]{A2003} Asmussen, Søren Applied probability and queues, Applications of Mathematics, 51, Springer, 2003, xii+438 pages \href{https://ahl.centre-mersenne.org/item/AHL_2019__2__281_0/#A2003}{Publisher reference}.
\bibitem[AZ09]{AZ09} Alexander, Kenneth S.; Zygouras, Nikos Quenched and annealed critical points in polymer pinning models, Commun. Math. Phys., Volume 291 (2009) no. 3, pp. 659-689 \href{https://ahl.centre-mersenne.org/item/AHL_2019__2__281_0/#AZ09}{Publisher reference}.
\bibitem[BCP + 14]{BCPSZ14} Berger, Quentin; Caravenna, Francesco; Poisat, Julien; Sun, Rongfeng; Zygouras, Nikos The critical curve of the random pinning and copolymer models at weak coupling, Commun. Math. Phys., Volume 326 (2014) no. 2, pp. 507-530 \href{https://ahl.centre-mersenne.org/item/AHL_2019__2__281_0/#BCPSZ14}{Publisher reference}.
\bibitem[Ber13]{B13} Berger, Quentin Comments on the influence of disorder for pinning model in correlated Gaussian environment, ALEA, Lat. Am. J. Probab. Math. Stat., Volume 10 (2013) no. 2, pp. 953-977 \href{https://ahl.centre-mersenne.org/item/AHL_2019__2__281_0/#B13}{Publisher reference}.
\bibitem[Ber14]{B14} Berger, Quentin Pinning model in random correlated environment: appearance of an infinite disorder regime, J. Stat. Phys., Volume 155 (2014) no. 3, pp. 544-570 \href{https://ahl.centre-mersenne.org/item/AHL_2019__2__281_0/#B14}{Publisher reference}.
\bibitem[BGT89]{BGT89} Bingham, Nicholas H.; Goldie, Charles M.; Teugels, Jozef L. Regular variation, Encyclopedia of Mathematics and Its Applications, 27, Cambridge University Press, 1989, 512 pages \href{https://ahl.centre-mersenne.org/item/AHL_2019__2__281_0/#BGT89}{Publisher reference}.
\bibitem[BL12]{BL12} Berger, Quentin; Lacoin, Hubert Sharp critical behavior for pinning models in a random correlated environment, Stochastic Processes Appl., Volume 122 (2012) no. 4, pp. 1397-1436 \href{https://ahl.centre-mersenne.org/item/AHL_2019__2__281_0/#BL12}{Publisher reference}.
\bibitem[BL16]{BL15} Berger, Quentin; Lacoin, Hubert Pinning on a defect line: characterization of marginal disorder relevance and sharp asymptotics for the critical point shift, J. Inst. Math. Jussieu, Volume 17 (2016) no. 2, pp. 205-346 \href{https://ahl.centre-mersenne.org/item/AHL_2019__2__281_0/#BL15}{Publisher reference}.
\bibitem[BP15]{BP15} Berger, Quentin; Poisat, Julien On the critical curve of the pinning and copolymer models in correlated Gaussian environment, Electron. J. Probab., Volume 20 (2015), 71, 35 pages \href{https://ahl.centre-mersenne.org/item/AHL_2019__2__281_0/#BP15}{Publisher reference}.
\bibitem[CdH13a]{CadH13} Caravenna, Francesco; den Hollander, Frank A general smoothing inequality for disordered polymers, Electron. Commun. Probab., Volume 18 (2013), 76, 15 pages \href{https://ahl.centre-mersenne.org/item/AHL_2019__2__281_0/#CadH13}{Publisher reference}.
\bibitem[CdH13b]{CdH13} Cheliotis, Dimitris; den Hollander, Frank Variational characterization of the critical curve for pinning of random polymers, Ann. Probab., Volume 41 (2013) no. 3B, pp. 1767-1805 \href{https://ahl.centre-mersenne.org/item/AHL_2019__2__281_0/#CdH13}{Publisher reference}.
\bibitem[CSZ16]{CSZ16} Caravenna, Francesco; Sun, Rongfeng; Zygouras, Nikos The continuum disordered pinning model, Probab. Theory Relat. Fields, Volume 164 (2016) no. 1-2, pp. 17-59 \href{https://ahl.centre-mersenne.org/item/AHL_2019__2__281_0/#CSZ16}{Publisher reference}.
\bibitem[CSZ17]{CSZ16+} Caravenna, Francesco; Sun, Rongfeng; Zygouras, Nikos Polynomial chaos and scaling limits of disordered systems, J. Eur. Math. Soc., Volume 19 (2017) no. 1, pp. 1-65 \href{https://ahl.centre-mersenne.org/item/AHL_2019__2__281_0/#CSZ16+}{Publisher reference}.
\bibitem[CTT17]{CTT16+} Caravenna, Francesco; Toninelli, Fabio Lucio; Torri, Niccoló Universality for the pinning model in the weak coupling regime, Ann. Probab., Volume 45 (2017) no. 4, pp. 2154-2209 \href{https://ahl.centre-mersenne.org/item/AHL_2019__2__281_0/#CTT16+}{Publisher reference}.
\bibitem[DGLT09]{DGLT09} Derrida, Bernard; Giacomin, Giambattista; Lacoin, Hubert; Toninelli, Fabio Lucio Fractional moment bounds and disorder relevance for pinning models, Commun. Math. Phys., Volume 287 (2009) no. 3, pp. 867-887 \href{https://ahl.centre-mersenne.org/item/AHL_2019__2__281_0/#DGLT09}{Publisher reference}.
\bibitem[dH09]{dH09} den Hollander, Frank Random Polymer models: École d’Été de Probabilités de Saint-Flour XXXVII-2007, Lecture Notes in Mathematics, 1974, Springer, 2009, xiii+258 pages \href{https://ahl.centre-mersenne.org/item/AHL_2019__2__281_0/#dH09}{Publisher reference}.
\bibitem[Dur10]{Du10} Durrett, Rick Probability: Theory and examples, Cambridge Series in Statistical and Probabilistic Mathematics, 31, Cambridge University Press, 2010 \href{https://ahl.centre-mersenne.org/item/AHL_2019__2__281_0/#Du10}{Publisher reference}.
\bibitem[Fel68]{Feller68} Feller, William An introduction to probability and its applications. Vol. 1, John Wiley \& Sons, 1968, xviii+509 pages \href{https://ahl.centre-mersenne.org/item/AHL_2019__2__281_0/#Feller68}{Publisher reference}.
\bibitem[Fel71]{Feller71} Feller, William An introduction to probability and its applications. Vol. 2, John Wiley \& Sons, 1971, xxiv+669 pages \href{https://ahl.centre-mersenne.org/item/AHL_2019__2__281_0/#Feller71}{Publisher reference}.
\bibitem[Fre82]{Fr82} Frenk, Johannes B. G. The behaviour of the renewal sequence in case the tail of the waiting-time distribution is regularly varying with index -1, Adv. Appl. Probab., Volume 14 (1982) no. 4, pp. 870-874 \href{https://ahl.centre-mersenne.org/item/AHL_2019__2__281_0/#Fr82}{Publisher reference}.
\bibitem[GHM99]{GHM1999} Georgii, Hans-Otto; Häggström, Olle; Maes, Christian The random geometry of equilibrium phases (Phase transitions and critical phenomena), Volume 18, Academic Press Inc., 1999, pp. 1-142 \href{https://ahl.centre-mersenne.org/item/AHL_2019__2__281_0/#GHM1999}{Publisher reference}.
\bibitem[Gia07]{Gi2007} Giacomin, Giambattista Random polymer models, World Scientific, 2007, xvi+242 pages \href{https://ahl.centre-mersenne.org/item/AHL_2019__2__281_0/#Gi2007}{Publisher reference}.
\bibitem[Gia11]{Gi2011} Giacomin, Giambattista Disorder and Critical Phenomena Through Basic Probability Models. École d’Été de Probabilités de Saint-Flour XL-2010, Lecture Notes in Mathematics, 2025, Springer, 2011, xi+130 pages \href{https://ahl.centre-mersenne.org/item/AHL_2019__2__281_0/#Gi2011}{Publisher reference}.
\bibitem[GL62]{GL63} Garsia, Adriano; Lamperti, John A discrete renewal theorem with infinite mean, Comment. Math. Helv., Volume 37 (1962), pp. 221-234 \href{https://ahl.centre-mersenne.org/item/AHL_2019__2__281_0/#GL63}{Publisher reference}.
\bibitem[GT06]{GiT2006} Giacomin, Giambattista; Toninelli, Fabio Lucio Smoothing effect of quenched disorder on polymer depinning transitions, Commun. Math. Phys., Volume 266 (2006) no. 1, pp. 1-16 \href{https://ahl.centre-mersenne.org/item/AHL_2019__2__281_0/#GiT2006}{Publisher reference}.
\bibitem[GTL10]{GTL09} Giacomin, Giambattista; Toninelli, Fabio Lucio; Lacoin, Hubert Marginal relevance of disorder for pinning models, Commun. Pure Appl. Math., Volume 63 (2010) no. 2, pp. 233-265 \href{https://ahl.centre-mersenne.org/item/AHL_2019__2__281_0/#GTL09}{Publisher reference}.
\bibitem[Har74]{H74} Harris, A. B. Effect of random defects on the critical behaviour of Ising model, J. Phys. C, Solid State Phys., Volume 7 (1974) no. 9, pp. 1671-1692 \href{https://ahl.centre-mersenne.org/item/AHL_2019__2__281_0/#H74}{Publisher reference}.
\bibitem[JPS06]{JPS06} Jeon, Jae-Hyung; Park, Pyeong Jun; Sung, Wokyung The effect of sequence correlation on bubble statistics in double-stranded DNA, J. Chem. Phys., Volume 125 (2006), 164901 \href{https://ahl.centre-mersenne.org/item/AHL_2019__2__281_0/#JPS06}{Publisher reference}.
\bibitem[Lac10]{L10} Lacoin, Hubert The martingale approach to disorder irrelevance for pinning models, Electron. Commun. Probab., Volume 15 (2010), pp. 418-427 \href{https://ahl.centre-mersenne.org/item/AHL_2019__2__281_0/#L10}{Publisher reference}.
\bibitem[LS17]{LS16+} Lacoin, Hubert; Sohier, Julien Disorder relevance without Harris Criterion: the case of pinning model with \(\gamma\)-stable environment, Electron. J. Probab., Volume 22 (2017), 50, 26 pages \href{https://ahl.centre-mersenne.org/item/AHL_2019__2__281_0/#LS16+}{Publisher reference}.
\bibitem[Nag12]{Na2012} Nagaev, Sergey V. The renewal theorem in the absence of power moments, Theory Probab. Appl., Volume 56 (2012) no. 1, pp. 166-175 \href{https://ahl.centre-mersenne.org/item/AHL_2019__2__281_0/#Na2012}{Publisher reference}.
\bibitem[Poi13a]{P13a} Poisat, Julien On quenched and annealed critical curves of random pinning model with finite range correlations, Ann. Inst. Henri Poincaré, Probab. Stat., Volume 49 (2013), pp. 456-482 \href{https://ahl.centre-mersenne.org/item/AHL_2019__2__281_0/#P13a}{Publisher reference}.
\bibitem[Poi13b]{P13} Poisat, Julien Ruelle-Perron-Frobenius operator approach to the annealed pinning model with Gaussian long-range correlated disorder, Markov Process. Relat. Fields, Volume 19 (2013) no. 3, pp. 577-606 \href{https://ahl.centre-mersenne.org/item/AHL_2019__2__281_0/#P13}{Publisher reference}.
\bibitem[Spi13]{Spitzer} Spitzer, Frank Principles of random walk, Graduate Texts in Mathematics, 34, Springer, 2013 \href{https://ahl.centre-mersenne.org/item/AHL_2019__2__281_0/#Spitzer}{Publisher reference}.
\bibitem[ST94]{ST94} Samorodnitsky, Gennady; Taqqu, Murad S. Stable non-Gaussian random processes, Stochastic Modeling, Chapman \& Hall, 1994, xviii+632 pages \href{https://ahl.centre-mersenne.org/item/AHL_2019__2__281_0/#ST94}{Publisher reference}.
\bibitem[Ton08]{T08} Toninelli, Fabio Lucio A replica-coupling approach to disordered pinning models, Commun. Math. Phys., Volume 280 (2008) no. 2, pp. 389-401 \href{https://ahl.centre-mersenne.org/item/AHL_2019__2__281_0/#T08}{Publisher reference}.
\bibitem[WH83]{WH83} Weinrib, Abel; Halperin, Bertrand I. Critical phenomena in systems with long-range-correlated disorder, Phys. Rev. B, Volume 27 (1983), pp. 413-427 \href{https://ahl.centre-mersenne.org/item/AHL_2019__2__281_0/#WH83}{Publisher reference}.
\end{thebibliography}
\end{document}
